%%%PAGE 259 rot=0 legibility=good summary=人撑在两面垂直的墙受力图；x-t图像题（甲乙两车，粘贴题）与解答；简易测a；匀圆不用动能定理；斜抛结论；多物体受力分析要点
%%%BLOCK id=259-01 ch=02 sec=受力分析·摩擦力 kind=diagram alt=none cont=none y=0.00-0.15
\textbf{人撑在两面垂直的墙}
%FIG p259_a 0.07 0.00 0.25 0.15
\notefig[0.3]{figures/p259_a.png}
%%%END
%%%BLOCK id=259-02 ch=01 sec=运动图像·x-t图像 kind=problem alt=none cont=none y=0.14-0.43
\begin{problem}
甲、乙两车在同一条直道上行驶，它们的位置坐标 $x$ 随时间 $t$ 变化的关系如图所示。已知乙车做匀变速直线运动，其图线最低点为 $P$，则下列说法正确的是（\quad）
% 原稿：粘贴的印刷题；「位置坐标 x」被手画椭圆圈出；括号处画了一个圈，并有一条长弧线连到选项D（示意答案为D）
%FIG p259_b 0.765 0.130 0.985 0.270
\notefig[0.3]{figures/p259_b.png}
\par A.\ \struck{甲车的初速度为$10\unit{m/s}$}
\par B.\ \struck{两车第一次相遇时运动方向相同}
\par C.\ 乙车的加速度大小\struck{为$1.2\unit{m/s^2}$}
\par D.\ 乙车的初始位置坐标 $x_0=90\unit{m}$
\begin{solution}
\[ x=a(t-10)^2+10 \]
代入 $(5,30)$\par
$a=\cdots$\par
$x_0=90$
\end{solution}
\end{problem}
% 题下圈起来的字
\fbox{坑}
%%%END
%%%BLOCK id=259-03 ch=03 sec=牛顿第二定律·简易加速度计 kind=diagram alt=none cont=none y=0.36-0.50
\textbf{简易测 $a$}
%FIG p259_c 0.05 0.36 0.25 0.50
\notefig[0.25]{figures/p259_c.png}
%%%END
%%%BLOCK id=259-04 ch=04 sec=圆周运动·匀速圆周 kind=note alt=06 cont=none y=0.53-0.56
匀圆不用动能定理
%%%END
%%%BLOCK id=259-05 ch=04 sec=平抛·斜抛 kind=knowledge alt=none cont=none y=0.60-0.77
\textbf{斜抛结论}
%FIG p259_d 0.075 0.617 0.230 0.770
\notefig[0.25]{figures/p259_d.png}
斜抛在 $\alpha$ 角平分线处\ 水平距离最大
%%%END
%%%BLOCK id=259-06 ch=02 sec=受力分析·多物体 kind=method alt=03 cont=none y=0.76-0.83
多物体复杂受力分析\quad 质量分块、受力均摊，分析节点（重力可能不是 $mg$，是\unclear{等效}结果 % 此行写到页边被截断，括号未闭合
\par 几何、三角函数关系 % 原稿写在上一行「分析节点」上方
%%%END
%%%PAGEEND 259
